% ============================================================================
% 2000 年春季高考（京、皖）文科 第 21 题 图一：直角梯形 ABCD
% ----------------------------------------------------------------------------
% 来源：.work/phase4-b20/source/2000春季京皖文.pdf
%       SHA-256 b9552b8354735296f6ea1ed051085e094e4af0c14ff97ada2dbd4e579b14287f
%       第 1 页，内容流对象 47（图注「图一」为排版文字，不在本图内）
% 重绘方式：qpdf --qdf 取内容流 → 解析 m/l/c/re/S/f、w、d → 本地坐标
%          本地坐标 = 页面坐标 − (874.160, 359.960)，原点取顶点 A（梯形左下角）；
%          1 单位 = 1 TeX pt；页面为 A3（1190.55 × 841.89 pt）。
% 元素计数：实线路径 4（AB、BC、CD、DA），虚线路径 1（AC），标签 4（A, B, C, D）
% 线宽：全部 0.40 pt；虚线 on/off 各 2.99 pt
% 坐标基准与源位的偏差：端点四舍五入到 0.001 pt；全部为直线段，无简化。
% 标签：位置 = pdftotext -bbox 文本中心 − 基准；D 源位 (−5.070,37.990)，左移 1.6 pt；C 源位 (38.580,37.990)，右移 1.0 pt、上移 2.2 pt；A、B 未动。净距均 ≥ 2 pt
% ============================================================================
\begin{tikzpicture}[x=1pt, y=1pt,
    every node/.style={inner sep=0pt, font=\fontsize{10.5pt}{12.6pt}\selectfont}]
  % --- 实线：AB、BC、CD、DA（内容流顺序为 AB→BC→CD→DA）
  \draw[line width=0.4pt] (0,0) -- (68.040,0);
  \draw[line width=0.4pt] (68.040,0) -- (34.020,34.020);
  \draw[line width=0.4pt] (34.020,34.020) -- (0,34.020);
  \draw[line width=0.4pt] (0,34.020) -- (0,0);
  % --- 虚线：AC（划长 2.99 pt，内容流 [2.98883 2.98883] 0 d）
  \draw[line width=0.4pt, dash pattern=on 2.99pt off 2.99pt] (0,0) -- (34.020,34.020);
  % --- 标签
  \node at (-6.670,37.990) {$D$};
  \node at (39.580,40.190) {$C$};
  \node at (-4.930,-5.890) {$A$};
  \node at (72.710,-5.890) {$B$};
  % 图注（源页图下居中；主会话添加）
  \node at (34.020,-12.000) {图一};
\end{tikzpicture}
% ----------------------------------------------------------------------------
% 自验：A=(0,0)、B=(68.040,0)、C=(34.020,34.020)、D=(0,34.020)，与内容流 874.16/942.20/908.18/874.16、359.96/393.98 逐点吻合；AD=DC=34.020、AB=68.040，满足题设 AD=DC=½AB；标签：A、B 取源位；D 左移 1.6 pt、C 右移 1.0 pt 并上移 2.2 pt。净距 A 3.32、B 3.22、D 2.44、C 2.18 pt。探针编译：0 错误、0 警告。
% ============================================================================
